\documentclass[10pt,a4paper]{amsart}
\usepackage[margin=19mm]{geometry}
\usepackage{amsmath,amssymb,mathtools}
\usepackage[scr=rsfs,cal=euler]{mathalfa}
\usepackage{xfrac}
\usepackage[unicode,hidelinks,bookmarks=false]{hyperref}
\newcommand{\ie}{i.e.\ }
\newcommand{\cf}{cf.\ }
\newcommand{\resp}{respectively\ }
\newcommand{\Subsection}{\subsection}
\providecommand{\nobreakdash}{\nobreak\hbox{-}}
\setlength{\emergencystretch}{2em}
\multlinegap0pt
\usepackage{amssymb}
\usepackage[shortlabels]{enumitem}
\usepackage{mathtools}
\usepackage{xcolor}
\newtheorem{theorem}{Theorem}
\newcommand\FS{\mathsf{FS}}
\newcommand\Nat{\mathbf{N}}
\newcommand\RT{\mathsf{RT}}
\newcommand{\lex}{\mathrm{lex}}
\newcommand\sW{\mathrm{sW}}
\title{\texorpdfstring{Free sets, thin sets and rainbows for barriers}{Free sets, thin sets and rainbows for barriers}}
\author{Lorenzo Carlucci, Oriola Gjetaj}
\date{Source: arXiv:2605.03523v1 (CC BY 4.0)}
\begin{document}
\maketitle

\noindent\emph{Source and licence.} \texorpdfstring{Free sets, thin sets and rainbows for barriers}{Free sets, thin sets and rainbows for barriers}. Version arXiv:2605.03523v1 (CC BY 4.0), distributed by arXiv under the Creative Commons Attribution 4.0 International licence, http://creativecommons.org/licenses/by/4.0/, as recorded in the arXiv licence metadata for this exact version and shown on the abstract page https://arxiv.org/abs/2605.03523v1. Full licence notice: \texttt{licenses/arxiv-2605.03523-CC-BY-4.0.txt}.\par\smallskip

\noindent\textit{Editorial scope.} Counted unit: the article's theorem thm:fs\_red\_rt (source0.tex lines 298--300). The article's complete original proof of it is delivered with it (source0.tex lines 302--345, one \texttt{proof} environment of 44 lines). No mathematics is written, repaired or re-derived: the statement and the proof are the article's own lines, in the article's own notation, and no retained line is substituted or re-flowed.  No further source-local context is needed: every cross-reference of every retained line resolves inside the two delivered spans.  Nothing was omitted from any retained span, so the package carries no omission record.  No retained line contains a citation, so the wrapper carries no bibliography and no bibliography entry.  No retained line contains a figure environment, an image-inclusion command or drawing code, so no image is delivered and the package declares no asset dependency.  Editorial formatting only, in addition: an AMS \texttt{amsart} preamble that restates the article's own declarations and macros used by the retained lines, namely the environments \texttt{theorem} on separate counters, together with the standard packages those lines need.  The retained environments print in this document as Theorem 1 in order of appearance, whereas the article prints the same items with the numbering of its own full document.  The complete native source of the article as distributed in the arXiv e-print for this exact version is bundled unchanged as \texttt{sources/199793/source0.tex}.
\begin{theorem}\label{thm:fs_red_rt}
For each barrier $B$ with base $\Nat$, $\FS^B\leq_{\sW} \RT^{B^+}_2$. 
\end{theorem}

\begin{proof}
Let $B$ be a barrier with base $\Nat$ and let $f:B\to \Nat$ be an instance of $\FS^B$.
For $s = (s_0, \dots, s_n) \in B^+$, let $s\ominus 1 = (s_0 - 1, \dots, s_{n-1}-1)$. Then $s \ominus 1 \in B$. Define $g$ by recursion on $<_{\lex}$ as follows. 

\begin{equation*}
    g(s) = \begin{cases}
           0 \hspace{1.5cm} & \text{if } f(s\ominus 1) = s_i-1 \textit{ for some } i  < n,  \\
           1 - g(s[f(s\ominus 1)+1]) & \text{if } f(s\ominus 1) < s_0 - 1, \text{or, for some } i < n-1,\\
           \hspace{1.2cm} & f(s\ominus 1) \in (s_i-1, s_{i+1}-1),\\
           0 \hspace{1.5cm} & \text{if } f(s\ominus 1) \in (s_{n-1}-1, s_n-1), \\
           1 \hspace{1.5cm} & \text{otherwise } (\textit{if } f(s\ominus 1) \geq s_n)   
           \end{cases} \quad
\end{equation*}

Note that $g$ is well-defined: if $f(s \ominus 1) < s_0 - 1$ then $f(s \ominus 1) + 1 < s_0 = \min(s)$, so that 
$s[f(s\ominus 1)+1]$ is well-defined, is in $B^+$ and is $<_\lex$-smaller than $s$. If $f(s\ominus 1) \in (s_i-1, s_{i+1}-1)$ for some  $i < n - 1$, then $f(s\ominus 1) +1 \in (s_i, s_{i+1})$ for some $i< n-1$ and $s[f(s\ominus 1)+1]$ is defined, is in $B^+$ and is $<_\lex$-smaller than $s$. 
Let $H = \{x_0, x_1, \dots \}\subseteq \Nat^+$ be infinite such that $B^+ | H$ is $g$-monochromatic as given by $\RT^{B^+}_2$. We claim that $B\vert H^-$ is $f$-free, where $H^- = \{x_0 - 1, x_1 - 1, \dots \}$. 
Consider some $(s_0 - 1, \dots, s_{n - 1} - 1) \in B\vert H^-$. Then $\{s_0, \dots, s_{n - 1}\} \subseteq H$ and for every $k \in H/s_{n-1}$, $(s_0, \dots, s_{n-1}, k) \in B^+\vert H$. 
To show that $B\vert H^-$ is $f$-free it is sufficient to show that for some $k$, writing $s=(s_0,\ldots,s_{n-1}, k) \in B^+\vert H$, we have that if $f(s\ominus 1) \in H^-$ then $f(s\ominus 1) \in s\ominus 1$, since all elements of $B\vert H^-$ are of the form $s\ominus 1$ for some $s \in B^+\vert H$. 
There are two cases:

\textbf{Case 1:} $B^+\vert H$ is monochromatic for the color $0$. Take $s_n$ to be the next element of $H$ after $s_{n - 1}$ and write $s = (s_0, \dots, s_{n})$, then $g(s) = 0$. We consider three subcases:

\textbf{Subcase 1.1:} $f(s \ominus 1) = s_i - 1$ for some $i < n$. 
Then $f(s\ominus 1)\in s\ominus 1$, so we are done.

\textbf{Subcase 1.2:} $f(s \ominus 1) \in (s_{n - 1} - 1, s_{n} - 1)$. 
Then $f(s\ominus 1)\notin H^-$ by definition of $s_n$ as the next element of $H$ after $s_{n-1}$.

\textbf{Subcase 1.3:} $f(s \ominus 1) < s_0 - 1$ or  $f(s \ominus 1) \in (s_i - 1, s_{i+1} - 1)$ for some $i < s_0 - 1$ and $g(s[f(s \ominus 1)+1]) = 1$. 
Since $f(s\ominus 1)+1 \in H$, we have $s[f(s\ominus 1)+1] \in B^+\vert H$. But then 
$g(s[f(s \ominus 1)+1]) = 1$ contradicts the fact that $B^+\vert H$ is $g$-monochromatic for the color $0$. 

\textbf{Case 2:} $B^+\vert H$ is monochromatic for the color $1$. Take $s_n \in H$ bigger than $s_{n - 1}$ and write $s = (s_0, \dots, s_{n})$, then $s \in B^+\vert H$ and $g(s) = 1$. We consider two subcases:

\textbf{Subcase 2.1:} $f(s \ominus 1) \geq s_{n}$. This case is impossible, indeed, as $H$ is infinite, there exists some element $x \in H$ such that $x > f(s \ominus 1) + 1$ and therefore $f(s \ominus 1) \in (s_{n} - 1, x - 1)$, which leads to $g(s_0, \dots, s_{n - 1}, x) = 0$ contradicting the fact that $B^+\vert H$ is $g$-monochromatic for the color $1$. Note that by definition of $B^+$, $(s_0, \dots, s_{n - 1}, x)\in B^+$. 

\textbf{Subcase 2.2:} $f(s \ominus 1) < s_0 - 1$ or $f(s \ominus 1) \in (s_i - 1, s_{i+1} - 1)$ for some $i < n - 1$ and $g(s[f(s \ominus 1)+1]) = 0$.  
Since $f(s\ominus 1)+1 \in H$, we have that $s[f(s\ominus )+1] \in B^+\vert H$. But 
$g(s[f(s \ominus 1)+1]) = 0$ contradicts the fact that $B^+\vert H$ is $g$-monochromatic for the color $1$. 

Therefore $B\vert H^-$ is $f$-free.
Notice that $g$ is uniformly computable in $f$ and that $H^-$ is uniformly computable in $H$. Hence, $\FS^B \leq_{\sW} \RT^{B^+}_2$.
 \end{proof}

\end{document}
